\subsection{Using the sampling toolkit with diffusion models}
\label{subsec:diffusion-sampling-viewpoints}

The previous subsection answered the model-building question.
It identified the exact reverse kernels and showed how, in Euclidean models, the corresponding reverse-time quantities can be learned by regression.
If those exact reverse kernels were available, then by \Cref{thm:distributional-posterior-learning,prop:reverse-kernels-transport-general} iterating them would transport \(p_1\) back to \(p_0\).
Likewise, if the exact score \(s(t,x)=\nabla\log p_t(x)\) and velocity \(b(t,x)\) were known and the reverse dynamics from \Cref{prop:si-family-of-sdes} could be simulated accurately, then plain reverse simulation would already solve the unconditional sampling problem.

That statement is theoretically important, but it does not settle the practical sampling question.
It answers an idealized problem:
how do we generate one unconditional sample from \(p_0\) if exact reverse-time quantities are available?
In practice we work with learned approximations, finite time grids, and tasks that go beyond unconditional i.i.d.\ sampling.
So once a reverse bridge is available, the question is no longer ``what is diffusion?'' but ``which sampling problem are we solving on top of that bridge?''

No new sampling framework is needed.
The earlier ideas from the course reappear with the same notation.
Fix a decreasing grid \(1=t_K>\cdots>t_0=0\), and write
\[
  X_{t_{k-1}}^{\mathrm{prop}}
  \sim
  \hat K_{t_{k-1}\mid t_k}(\,\cdot\,\mid X_{t_k})
\]
for the approximate reverse kernel produced by a learned finite-step sampler such as \Cref{alg:bridge-sampling}.
This is meant to approximate the exact reverse kernel \(K_{t_{k-1}\mid t_k}\), not the endpoint posterior \(q_{0\mid t_k}\).
Accordingly, \(X_{t_{k-1}}^{\mathrm{prop}}\) is a proposal for the state at the earlier time \(t_{k-1}\), not a proposal for \(X_0\) unless \(k=1\).
The decreasing-noise family \(\{p_t\}_{t\in[0,1]}\) plays the same structural role as the target ladder in \Cref{subsec:tempering}:
instead of attacking the difficult endpoint law \(p_0\) in one step, the algorithm moves through a sequence of easier intermediate laws.
Once such a bridge is on the table, modern diffusion samplers mostly modify one of three things: the proposal family, the target family, or the state space.
We may want the same target law \(p_0\), but at lower computational cost.
We may want a conditional target law \(p_0(\cdot\mid y)\) rather than the unconditional one.
Or we may want a finite set of samples whose empirical distribution approximates the target law better than i.i.d.\ sampling would at the same budget.
In the notation of the course, these are three different ways to modify the sampling problem:
for efficient sampling, the target family \(\{p_t\}\) stays fixed and we redesign the proposal or add a local correction;
for guidance, the target family changes from \(\{p_t\}\) to the conditional family
\[
  p_t^y(dx)\coloneqq p_t(dx\mid y);
\]
for finite-particle approximation, the state space changes from \(\R^d\) to the product space \((\R^d)^N\) and the target changes from the product law \((p_t^y)^{\otimes N}\) to a coupled law on that product space.
More concretely, the first task points back to gradient-based MCMC in \Cref{subsubsec:langevin-diffusion,subsubsec:mala}, the second to conditioning by likelihood reweighting in \Cref{ex:conditioning-likelihood-reweighting} and sequential reweighting in \Cref{subsec:smc}, and the third to the empirical-measure viewpoint of particle methods in \Cref{subsubsec:smc-particles,subsubsec:smc-resampling}.

\subsubsection{Efficient sampling for the same target law}
\label{subsubsec:diffusion-efficient-sampling}

The first task keeps the target law \(p_0\) fixed and asks for a cheaper sampler.
Correctness is not the issue here; computation is.
With a learned model, the practical endpoint error comes from both approximation error in the learned field and discretization error along the time grid.
If controlling those errors requires a very fine grid, then reverse sampling becomes slow.
Because each reverse step requires one or more evaluations of the learned score or drift network, the natural cost measure is the number of network function evaluations (NFEs).
So ``more efficient sampling'' means lower endpoint error at a fixed NFE budget, or comparable endpoint error using fewer NFEs.

The important point is that efficiency does not mean changing the target.
We still want samples from the same law \(p_0\).
We are only changing how we approximate the reverse path.
This is exactly the distinction that appeared earlier in gradient-based MCMC:
keep the target fixed, then redesign the proposal and, if useful, locally correct it.

Two familiar levers now reappear.
The first is to treat \(\hat K_{t_{k-1}\mid t_k}\) as a predictor and then follow it with a short correction targeted at the new time level.
Formally, one replaces the bare proposal
\[
  X_{t_{k-1}}^{\mathrm{prop}}
  \sim
  \hat K_{t_{k-1}\mid t_k}(\,\cdot\,\mid X_{t_k})
\]
by the composite update
\[
  X_{t_{k-1}}
  \sim
  C_k(\,\cdot\,\mid X_{t_{k-1}}^{\mathrm{prop}}),
  \qquad
  X_{t_{k-1}}^{\mathrm{prop}}
  \sim
  \hat K_{t_{k-1}\mid t_k}(\,\cdot\,\mid X_{t_k}),
\]
where \(C_k\) is a short Markov kernel chosen so that its invariant law is, at least approximately, the current marginal \(p_{t_{k-1}}\).
The simplest example is a frozen-time Langevin corrector:
\begin{equation}
  X^{(\ell+1)}
  =
  X^{(\ell)}
  +
  \eta_k\, s(t_{k-1},X^{(\ell)})
  +
  \sqrt{2\eta_k}\,\xi_\ell,
  \qquad
  \xi_\ell\sim\mathcal{N}(0,I_d),
  \label{eq:diffusion-corrector-ula}
\end{equation}
initialized at \(X^{(0)}=X_{t_{k-1}}^{\mathrm{prop}}\).
If \(p_t\equiv \pi\) were fixed in time, then \eqref{eq:diffusion-corrector-ula} would be exactly the ULA update from \Cref{subsubsec:langevin-diffusion}.
Here the target changes with \(k\), so the same Langevin idea is used only as a short local repair step at the frozen time \(t_{k-1}\).
If density ratios were tractable, one could similarly add an accept/reject correction as in \Cref{subsubsec:mala}.
Such a corrector costs extra work locally, but it can still be efficient globally if it permits a coarser grid, fewer reverse steps, or a noticeably smaller endpoint error at the same total NFE budget.
This is the same proposal-versus-correction tradeoff as in \Cref{subsubsec:langevin-diffusion,subsubsec:mala}, now applied along the reverse bridge.

The second lever is to keep the same continuous-time reverse dynamics but change the proposal itself by using a better discretization.
When \(\varepsilon\equiv 0\), \Cref{prop:si-family-of-sdes} gives the probability-flow ODE
\begin{equation}
  \dot x_t=b(t,x_t),
  \label{eq:diffusion-pf-ode-local}
\end{equation}
whose marginals are again \(\{p_t\}\).
The most basic discretization is Euler:
\[
  x_{t_{k-1}}
  \approx
  x_{t_k} + (t_{k-1}-t_k)\,b(t_k,x_{t_k}).
\]
A more accurate two-stage rule first predicts with Euler and then averages the old and predicted drifts:
\[
  x^{\mathrm{prop}}
  =
  x_{t_k} + (t_{k-1}-t_k)\,b(t_k,x_{t_k}),
\]
\[
  x_{t_{k-1}}
  \approx
  x_{t_k}
  +
  \frac{t_{k-1}-t_k}{2}
  \Bigl(
    b(t_k,x_{t_k})
    +
    b(t_{k-1},x^{\mathrm{prop}})
  \Bigr).
\]
Both updates aim at the same path \(\{p_t\}\).
They differ only in the tradeoff between computation and discretization bias.
So the fast-sampling question is again a proposal-design question:
how much endpoint accuracy do we get from each unit of computation?
In diffusion language, proposal-plus-correction appears as predictor--corrector sampling \citep[Alg.~1]{song_sohl_dickstein_kingma_kumar_ermon_poole_2021_score_sde}.
Low-NFE deterministic samplers such as DDIM \citep[\S5]{song_meng_ermon_2021_ddim} and higher-order ODE solvers such as DPM-Solver \citep[\S3 and \S5.1]{lu_zhou_bao_chen_li_zhu_2022_dpm_solver} fit the same template:
they redesign \(\hat K_{t_{k-1}\mid t_k}\) so that the same endpoint task can be solved with fewer reverse evaluations.

\subsubsection{Guidance as posterior tilting}
\label{subsubsec:diffusion-conditional-sampling}

The second task changes the target law.
We may want a sample that matches external information \(y\): a class label, a text description, or a measurement.
Then the target is not \(p_0\), but the conditional law \(p_0(\cdot\mid y)\).
An unconditional reverse sampler can be perfectly valid for \(p_0\) and still be solving the wrong problem.

Guidance is the name used in the diffusion literature for this target change.
In the notation of these notes, the unconditional path \(\{p_t\}\) is replaced by the conditional path \(\{p_t^y\}\), where
\[
  p_t^y(dx)
  \coloneqq
  p_t(dx\mid y)
  =
  \frac{p_t(y\mid x)\,p_t(dx)}{\int p_t(y\mid x')\,p_t(dx')}.
\]
This is exactly the conditioning-by-likelihood-reweighting identity from \Cref{ex:conditioning-likelihood-reweighting}, now indexed by time \(t\).
Relative to importance sampling and SMC, the only structural difference is where that likelihood factor enters:
there it appears through weights after propagation, whereas here it is inserted directly into the reverse-time field.

To express the conditional task in the same notation, let \(Y\) be an external variable and assume the joint law of \((X_t,Y)\) has a density.
Then Bayes's rule gives the conditional score.

\begin{proposition}[Conditional scores split into prior and likelihood terms]
\label{prop:diffusion-guided-score-identity}
For every fixed \(t\) and every \(y\) with \(p_t(y\mid x)>0\),
\begin{equation}
  \nabla_x \log p_t(x\mid y)
  =
  \nabla_x \log p_t(x)
  +
  \nabla_x \log p_t(y\mid x).
  \label{eq:diffusion-guided-score-identity}
\end{equation}
\end{proposition}
\begin{proof}
Bayes's rule gives
\[
  p_t(x\mid y)=\frac{p_t(x)p_t(y\mid x)}{p_t(y)}.
\]
Taking \(\nabla_x\log\) and using that \(p_t(y)\) does not depend on \(x\) yields \eqref{eq:diffusion-guided-score-identity}.%
\end{proof}

This identity gives the sampling rule immediately.
To sample from the conditional path, one replaces the unconditional score \(s(t,x)\) in \eqref{eq:si-forward-backward-drifts} by
\[
  s_y(t,x)
  \coloneqq
  \nabla_x \log p_t(x\mid y)
  =
  s(t,x)+\nabla_x \log p_t(y\mid x).
\]
So the reverse-time sampler is unchanged except for an added likelihood-gradient term.
In the language of these notes, guidance is posterior tilting along the path \(\{p_t\}\):
the prior score \(s(t,x)\) describes what the diffusion model prefers on its own, while the likelihood term \(\nabla_x\log p_t(y\mid x)\) pulls the sample toward states compatible with the side information \(y\).

Sometimes one intentionally strengthens that tilt and uses
\begin{equation}
  s_w(t,x;y)
  \coloneqq
  s(t,x)+w\,\nabla_x \log p_t(y\mid x),
  \qquad w\ge 0.
  \label{eq:diffusion-weighted-conditional-score}
\end{equation}
When \(w=1\), this is the exact conditional score.
When \(w>1\), it is a more concentrated target:
\[
  s_w(t,x;y)
  =
  \nabla_x \log \bigl(p_t(x)p_t(y\mid x)^w\bigr).
\]
So \(w>1\) does not leave the conditional law unchanged; it produces a tempered version that puts more mass on states that fit \(y\) especially well.
This is the same power-target idea as in \Cref{subsec:tempering}, but now it is applied only to the likelihood factor rather than to the whole target law.
When \(w=1\) we are solving the exact conditional sampling problem.
When \(w>1\) we are intentionally biasing the target toward stronger agreement with \(y\), typically to trade some faithfulness to the exact conditional law for sharper or more recognizable samples.

If a model supplies both an unconditional estimate \(s_{\varnothing}(t,x)\approx s(t,x)\) and a conditional estimate \(s_y(t,x)\approx \nabla_x\log p_t(x\mid y)\), then the common interpolation
\[
  s_{\mathrm{cfg}}(t,x;y)
  \coloneqq
  s_{\varnothing}(t,x)+w\bigl(s_y(t,x)-s_{\varnothing}(t,x)\bigr)
\]
reduces, in the exact case, to the weighted tilt \eqref{eq:diffusion-weighted-conditional-score}.
This is the form often used in practice.

The diffusion literature uses different names depending on how the extra likelihood information is obtained.
When the likelihood-gradient term is supplied by a separate noisy-time classifier, the method is called classifier guidance \citep[\S4]{dhariwal_nichol_2021_diffusion_models_beat_gans}.
When the model itself is trained to produce both unconditional and conditional outputs and the two are interpolated as above, the method is called classifier-free guidance \citep[\S2]{ho_salimans_2022_classifier_free_guidance}.
In our notation both methods are instances of the same posterior-tilting identity; they differ only in how \(\nabla_x\log p_t(y\mid x)\) is approximated.
The same formula also covers inverse problems once \(y\) is interpreted as a measurement, so we do not need a separate sampling principle for that case.

\subsubsection{A sample set as a distributional approximation}
\label{subsubsec:diffusion-diverse-sampling}

The third task changes the objective again.
Sometimes the output of interest is not one correct sample but a set of \(N\) samples that approximates the target law well as a whole.
Then the natural object is the empirical measure
\[
  \hat\mu_N
  \coloneqq
  \frac1N\sum_{i=1}^N \delta_{x_0^i},
\]
and the practical goal is that \(\hat\mu_N\) approximate the target law \(p_0(\cdot\mid y)\) well.
Equivalently, for many test functions \(f\), we want
\[
  \frac1N\sum_{i=1}^N f(x_0^i)
  \approx
  \E\bigl[f(X_0)\mid Y=y\bigr].
\]
This is the empirical-measure viewpoint from the particle methods earlier in the course, especially \Cref{subsubsec:smc-particles,subsubsec:smc-resampling}.
In particular, \Cref{prop:particle-cloud-unbiased-variance} shows that once the objective is empirical approximation, exact independence is sufficient but not necessary:
a coupled cloud can reduce empirical-average variance for important test functions and therefore outperform an i.i.d.\ cloud at the same particle budget.

To write the problem cleanly in the course notation, let
\[
  x_t^{1:N}\coloneqq (x_t^1,\dots,x_t^N)\in(\R^d)^N.
\]
If we insist on exact i.i.d.\ sampling, then the natural \(N\)-particle target is just the product law
\[
  (p_t^y)^{\otimes N}(dx_t^{1:N})
  =
  \prod_{i=1}^N p_t(dx_t^i\mid y).
\]
Independent sampling gives the correct one-particle marginals, but it does not by itself solve the finite-\(N\) approximation problem.
At finite \(N\), an i.i.d.\ batch can cluster and leave important regions underrepresented.
So if the goal is to approximate the law well with a fixed particle budget, it can be worthwhile to replace the product target by a coupled law:
\begin{equation}
  \Pi_t(dx_t^{1:N}\mid y)
  \propto
  (p_t^y)^{\otimes N}(dx_t^{1:N})\,
  \exp\!\bigl(\lambda R_t(x_t^{1:N})\bigr).
  \label{eq:diffusion-diverse-path-target}
\end{equation}
Here \(\lambda\ge 0\) controls the interaction strength, and \(R_t\) is chosen to reward spread, coverage, or repulsion.
When \(\lambda=0\), \eqref{eq:diffusion-diverse-path-target} reduces to the i.i.d.\ product law.
When \(\lambda>0\), the target is a coupled cloud chosen to improve the empirical approximation obtained from \(N\) particles.

Now the state variable is the whole configuration \(X_t^{1:N}\), not a single particle.
This is an augmented-state construction in exactly the sense of \Cref{sec:augmented-state-sampling}:
the state is enlarged from one particle to a whole interacting configuration so that local computations can produce a better global approximation.
Once the problem is written this way, the earlier tools apply directly:
one can run Langevin or Metropolis updates on the joint configuration, or propagate and reweight a particle population as in \Cref{subsec:smc}.
The interaction term is not a correction applied after sampling.
It is part of the target itself.

This is the viewpoint taken explicitly in Particle Guidance, which introduces a time-dependent joint-particle potential in order to obtain non-i.i.d.\ diverse diffusion samples \citep[\S3]{corso_xu_de_bortoli_barzilay_jaakkola_2023_particle_guidance}.
In the notation of these notes, that paper is an interacting-particle sampler on a product-space diffusion target.
The phrase ``non-i.i.d.'' there simply means that the joint law \(\Pi_t\) no longer factorizes into a product of one-particle marginals.

These three tasks explain much of the modern sampling literature around diffusion models.
Predictor--corrector sampling, DDIM, and DPM-Solver redesign the proposal or local correction while keeping the same target.
Classifier guidance and classifier-free guidance implement posterior tilting along the path.
Particle Guidance replaces the product target by an interacting law on the particle cloud.
The reverse bridge provides the path \(\{p_t\}\); the sampling toolkit from the earlier units does the rest.
